% ===================================================================
%  نقشہ نمبر 7 — سیال وچ گراں۔ — بہار وچ پیلیاں۔
%  Traduction 2026 d'apres texto/io, controlee sur texto/fr, texto/en
%  et texto/hi. Consigne : voir texto/pnb/10-naqsha-01.tex.
% ===================================================================

%%K t07-noto-1 noto
\VUnotes{143.2mm}{%
(*) کھاݨے دے ویروے لئی نقشے 6 تے 13 ویکھو۔%
}

%%K t07-apar-1 apar
\VUsaut{4.0mm}
\VUcentre{13.2pt}{90}{دوجا سلسلہ}
\VUsaut{3.0mm}
\VUfilet{16mm}
\VUsaut{5.6mm}
\VUtitre{15.0pt}{60}{نقشہ نمبر 7}
\VUsaut{3.0mm}

%%K t07-tit-1 sub
\VUcentre{12.0pt}{0}{\textit{پہلا منظر۔}}
\VUsaut{3.0mm}

%%K t07-tit-2 sub
\VUpk{10.2pt}{40}{سیال وچ گراں۔}
\VUsaut{4.0mm}

%%K t07-c1-01-1 p
1. --- نویں ورھے دیاں چھٹیاں وچ میں اپݨے پیو نال نواں گراں گیا، اک نِکا پنڈ
جِتھے اوہناں نے کجھ کم نبیڑنا سی۔

%%K t07-c1-02-1 p
2. --- اسیں اوہ \VUgras{ڈاک گڈی} \textsuperscript{(11)} لئی جیہڑی شہر تے اوس پنڈ
دے وچکار چلدی اے؛ پر پالا بہت سی، تے اِنّی اوکھی گڈی وچ سفر رتی وی سوہݨا نہ
رہیا۔

%%K t07-c1-03-1 p
3. --- ایسے لئی سانوں سچ مُچ خوشی ہوئی جدوں اسیں \VUgras{گڈی والے} نوں چونک وچ
\textit{چٹے رِچھ} دی \VUgras{سرائے} \textsuperscript{(2)} دے سامݨے گڈی روکدے
ویکھیا۔ \VUgras{سرائے والا} \textsuperscript{(4)} اپݨی دہلیز اُتے ساڈی اُڈیک کر
رہیا سی، چُپ چاپ اپݨی \VUgras{چِلم} پیندا ہویا۔

%%K t07-c1-03-2 p
اوہنے سانوں اندر سدیا، تے چولھے وچ چٹکدی ویل دی اگ دے سامݨے بہݨ لئی مینوں دوجی
واری کہاؤݨ دی لوڑ نہ پئی۔ فیر مینوں اوس کمبݨ لاؤݨ والی \VUgras{اُتری ہوا} دی
کوئی فکر نہ رہی جیہڑی پرانے \VUgras{بورڈ} \textsuperscript{(3)} نوں چرچراندی سی تے
چمنی توں نکلدے \VUgras{دھوئیں} \textsuperscript{(1)} نوں کالے وَروَلیاں وچ اُڈا لے
جاندی سی۔

%%K t07-c1-04-1 p
4. --- دوپہر دے کھاݨے دا ویلا سی۔ تے ہور دیر کیتے بغیر اسیں اوس سادے پر سوادی
کھاݨے نال پورا انصاف کیتا جیہڑا سانوں وَرتایا گیا (*)۔

%%K t07-tit-3 sub
\VUpk{10.2pt}{40}{کجھ ملاقاتاں۔}
\VUsaut{3.0mm}

%%K t07-c1-05-1 p
5. --- کھاݨے مگروں میں اپݨے پیو نال گیا، جیہڑے بیمے دے ایجنٹ نیں، اوہناں دے
گاہکاں نوں ملݨ۔ پہلاں اسیں \VUgras{لوہار} \textsuperscript{(5)} کول گئے، جیہڑا
سرائے دی بغل وچ رہندا اے۔ اوہ اک گھوڑے دا نعل ٹھوکݨ وچ لگا ہویا سی۔ میں اوہدے
مددگار دے زور دی داد دِتی، جیہڑا گھوڑے دا \VUgras{سُم} اپݨے چمڑے دے پرݨے اُتے
کس کے پھڑی ہویا سی، جدوں تیکر لوہار بھارے ہتھوڑے دیاں چوٹاں نال \VUgras{نعل}
\textsuperscript{(6)} گڈدا رہیا۔

%%K t07-c1-06-1 p
6. --- فیر اسیں گراں دے \VUgras{موچی} \textsuperscript{(7)} بُڈھے جیکب کول گئے۔
بھانویں اوہدے بورڈ اُتے اک شاندار \VUgras{بوٹ} بݨیا ہویا سی، اوہ پرانے گھِسے
جُتیاں وچ نویں تلے لا کے کم چلا رہیا سی۔

%%K t07-c1-06-2 p
بُڈھے نے ہسدے ہوئے آکھیا : « شہر وچ میں “جُتے بݨاؤݨ والا” سی تے میرے گاہک ای نہ
سن۔ ایتھے میں “موچی” واں تے کم دی کوئی گھاٹ نئیں۔ »

%%K t07-c1-07-1 p
7. --- اوہنے سانوں اپݨے گوانڈھی \VUgras{حجام} \textsuperscript{(8)} دی گل دسی،
جیدے گاہک سارے بے تسلے رہندے نیں۔ اوہدے \VUgras{اُسترے} ہمیشہ کھنڈے رہندے نیں تے
اوہدی \VUgras{قینچی} کدے ٹھیک نئیں کٹدی۔ پر کیوں جے اوہو ای گراں دا کلا حجام اے،
لوکاں کول اوسے توں حجامت کراؤݨ تے وال کٹواؤݨ توں بغیر چارہ نئیں۔ اوہ اُتوں کنجوس
تے رُکھا وی اے۔

%%K t07-c1-08-1 p
8. --- خوش قسمتی نال دوجے \VUgras{گوانڈھی} اوہدے جیہے نئیں۔ \VUgras{گڈی بݨاؤݨ
والا} \textsuperscript{(9)}، مثال لئی، بھلا بندہ اے، محنتی تے مددگار۔ اوہنوں گڈی
مرمت نوں دیݨا سُکھ دی گل اے۔ اوہدا کم ہمیشہ صاف سُتھرا ہوندا اے۔

%%K t07-c1-09-1 p
9. --- بُڈھے جیکب نوں \VUgras{زین ساز} \textsuperscript{(10)} توں وی کوئی شکایت نہ
سی، جیدی اوہ کم دے دباء وچ کدے کدے \VUgras{ساز} سیوݨ وچ مدد کردا اے۔

%%K t07-c1-09-2 p
میرے پیو نے اوہدے کولوں اوہدے ٹبر بارے پُچھیا۔ « سارے راضی باضی نیں۔ میری پوتری
\VUgras{سکول} \textsuperscript{(12)} وچ اے۔ اوہدی \VUgras{ماسٹرݨی}
\textsuperscript{(13)} اوہدے توں بہت خوش اے؛ اوہ چنگی \VUgras{طالبہ}
\textsuperscript{(14)} اے۔ اوہ میرے اوس شرارتی پُتر جیہی نئیں؛ اوہنوں کھیڈ توں
بغیر کجھ نئیں سُجھدا۔ \VUgras{مکھ ماسٹر} \textsuperscript{(15)} اوہنوں کجھ وی نئیں
سکھا سکدے۔ »

%%K t07-tit-4 sub
\VUsaut{3.0mm}
\VUpk{10.2pt}{40}{گراں دیاں گلاں۔}
\VUsaut{3.0mm}

%%K t07-c1-10-1 p
10. --- فیر بُڈھے نے سانوں گراں دیاں خبراں سُݨائیاں۔ اک نواں سکول بݨنا سی، کیوں
جے \VUgras{گراں دے گھر} والا بہت نِکا پے گیا سی۔ ایہہ سوکھا ہوندا، کیوں جے
\VUgras{چکی والے} دے باغ دا اک حصہ خرید لیݨا ای بہت سی۔ وچارے لئی ایہہ بڑا چنگا
رہندا۔ جدوں دا پالا پیا، \VUgras{چکی} \textsuperscript{(16)} دے \VUgras{پُڑ} کݨک
نہ پیہہ سکے، کیوں جے \VUgras{پہیا} \textsuperscript{(17)} \VUgras{برف}
\textsuperscript{(25)} وچ پھس گیا سی، جیہڑی \VUgras{نالے} \textsuperscript{(23)} دی
سی۔

%%K t07-c1-11-1 p
11. --- « بݨتر دی گل چلی تے — ساڈا نواں \VUgras{گرجا} \textsuperscript{(18)}
ویکھیا جے؟ برف دی چادر دے تھلے اوہ بڑا سوہݨا لگدا اے۔ \VUgras{کاں}
\textsuperscript{(19)}، جیہڑے اپݨا پرانا گھنٹہ گھر ہُݨ پچھاݨدے ای نئیں، اداس ہو کے
\VUgras{قبرستان} \textsuperscript{(20)} دے وڈے رُکھ اُتے بیٹھے رہندے نیں۔ »

%%K t07-c1-12-1 p
12. --- قبرستان دا ناں آندے ای موچی نوں اوہ لوک یاد آ گئے جیہناں نوں میرے پیو
جاݨدے سن تے جیہڑے پچھلے ورھے ٹُر گئے۔ « کِنّیاں \VUgras{قبراں}
\textsuperscript{(21)}، میرے بھلے جناب، \VUgras{سرو} \textsuperscript{(22)} دے تھلے
ہور جُڑ گئیاں! موت ساڈے بُڈھے منشی نوں لے گئی۔ سارا گراں اوہناں دے
\VUgras{جنازے} وچ سی۔ »

%%K t07-c1-13-1 p
13. --- « اوہ رہے اوہناں دے وارث، نالے اُتے برف تِلک رہے نیں۔ تھوڑی دیر پہلاں ای
آئے سن پئی میں اوہناں دی \VUgras{برف جُتی} \textsuperscript{(26)} دا تسمہ جوڑ
دیواں، جیہڑا ٹُٹ گیا سی۔ »

%%K t07-c1-14-1 p
14. --- « تے اوہ رہیا، نالے دے کنڈھے، نواں \VUgras{گراں دا چوکیدار}
\textsuperscript{(27)}۔ اوہ پرانا سپاہی اے تے گراں دے شرارتی مُنڈیاں دا کال۔ اوہدے
\VUgras{پٹے} \textsuperscript{(29)} تے اوہدی \VUgras{جھالر والی ٹوپی}
\textsuperscript{(28)} دِسدے ای اوہ نس کھڑدے نیں۔ »

%%K t07-c1-15-1 p
15. --- « اوہو! اوہ رہیا بُڈھا جوزف، نانبائی لئی \VUgras{لکڑاں دیاں گٹھڑیاں دی
گڈی} \textsuperscript{(30)} لیاندا ہویا۔ --- ٹھیک آکھیا، میرے پیو بولے، میں اوہدے
\VUgras{خچراں} \textsuperscript{(31)} دی جوڑی پچھاݨدا واں۔ مینوں اوہدے نال اک گل
کرنی اے، ایہو موقع اے۔ سلام، بُڈھے جیکب۔ »

%%K t07-c1-16-1 p
16. --- اسیں باہر نکل ای رہے سی پئی گراں دے نِکے سکول توں چھُٹ کے
\VUgras{برف دی تِلک پٹی} \textsuperscript{(33)} اُتے ٹُٹ پئے، ایس خطرے نال پئی
\VUgras{ڈِگݨ} \textsuperscript{(34)} اُتے ہتھ پیر تُڑا بہݨ۔ اوہناں وچوں اک نے گڈی
والے دے گھروں اپݨی \VUgras{سلیج} \textsuperscript{(32)} لیائی، اوہدے اُتے اپݨا اک
ساتھی بہایا، تے دوڑ پیا۔

%%K t07-c1-17-1 p
17. --- دوجے اک \VUgras{برف دے ڈھیر} \textsuperscript{(37)} ول لپکے، جیہڑا
\VUgras{پُل} \textsuperscript{(40)} کول سی، تے اوس \VUgras{برف دے بندے}
\textsuperscript{(39)} اُتے مارن لگے جیہڑا اوہناں کل بݨایا سی، جیہنوں اوہناں اک
ٹوپی، اک چِلم تے مونڈھے اُتے لٹکیا اک بستہ دِتا ہویا سی۔

%%K t07-c1-18-1 p
18. --- اوہناں نوں برفیلی ہوا تے پالے دی پرواہ نہ سی۔ اوہناں وچوں بہتیاں کول
\VUgras{گلوبند} \textsuperscript{(35)} تیکر نہ سی، تے \VUgras{برف دے گولیاں}
\textsuperscript{(38)} دی لڑائی مگروں نیمے ہوݨ لئی مُٹھ بھر تپدے
\VUgras{چیسٹ نٹ} \textsuperscript{(42)} ای اوہناں نوں بہت سن، جیہڑے اوہ بُڈھی جیکلین
\VUgras{ویچݨ والی} کولوں خریددے سن، جیہڑی اپݨی بہت پتلی \VUgras{چادر}
\textsuperscript{(43)} وچ پالے نال کمبدی رہندی اے۔

%%K t07-tit-5 sub
\VUsaut{3.0mm}
\VUcentre{12.0pt}{0}{\textit{دوجا منظر۔}}
\VUsaut{3.0mm}

%%K t07-tit-6 sub
\VUpk{10.2pt}{40}{بہار وچ پیلی دا گھر۔}
\VUsaut{4.0mm}

%%K t07-c2-01-1 p
1. --- سیال دے سارے سُکھاں دے باوجود \VUgras{بہار} دے مُڑ آوݨ دا اسیں ڈونگھے آنند
نال سواگت کرنے آں۔

%%K t07-c2-01-2 p
مئی دے مہینے وچ شام نوں پیلی دا وِہڑا کِہو جیہی جوش بھری تصویر بݨاندا اے!

%%K t07-c2-02-1 p
2. --- اُفق اُتے \VUgras{سورج} \textsuperscript{(1)} ہولی ہولی ڈُب رہیا اے، تے
\VUgras{پنڈ} \textsuperscript{(3)} نوں اپݨیاں لال \VUgras{کِرناں} \textsuperscript{(2)}
نال بھر رہیا اے۔ محنتی \VUgras{ہلوان} \textsuperscript{(6)} اپݨی \VUgras{پیلی}
\textsuperscript{(4)} وا ہُݨ دی کاہلی وچ اے۔

%%K t07-c2-02-2 p
اپݨے \VUgras{ہل} \textsuperscript{(5)} نال، جیہنوں اک تکڑا \VUgras{گھوڑا}
\textsuperscript{(7)} کھِچدا اے، اوہ \VUgras{سیاڑ} \textsuperscript{(8)} کھِچدا اے،
جیدے وچ \VUgras{بیج بیجݨ والا} \textsuperscript{(9)} مُٹھ مُٹھ اوہ \VUgras{بی}
چھَٹے گا جیدے توں سُنہری سِٹے نکلݨ گے۔

%%K t07-c2-03-1 p
3. --- اپݨیاں دِرانیاں نال \VUgras{گھاہ کٹݨ والیاں} \textsuperscript{(11)} نے
میدان دا گھاہ کٹ لیا اے، سُکاؤݨ والیاں نے \VUgras{سُکا گھاہ} \textsuperscript{(10)}
نوں \VUgras{پنجیاں} \textsuperscript{(12)} تے کھرپیاں نال پرت پرت کے سُکا لیا اے،
تے ہُݨ اوہ مُڑ رہے نیں۔

%%K t07-c2-03-2 p
کجھ بلداں نال کھِچی بھاری گڈی دے اگے اگے ٹُردے نیں، مونڈھے اُتے اپݨے اوزار چُکی۔
دوجے، ودھ سست، مہکدے سُکے گھاہ اُتے آرام نال پسر گئے نیں۔

%%K t07-c2-04-1 p
4. --- ٹُردے ٹُردے اوہ \VUgras{مالی} \textsuperscript{(16)} نوں یاری نال سلام کردے
نیں، جیہڑا \VUgras{باغ} \textsuperscript{(13)} وچ لگا اے، \VUgras{پھلاں والے
رُکھاں} دی چھانگ کردا تے \VUgras{ناشپاتی دے رُکھاں} \textsuperscript{(14)} وچ قلم
لاندا۔ \VUgras{آلو بخارے دے رُکھ} \textsuperscript{(15)} پھُلاں نال لدے نیں، تے چنگی
رُوڑی پائے \VUgras{ترکاری دے باغ} \textsuperscript{(17)} وچ ترکاریاں، گوبھی تے سلاد
پہلاں توں چنگیاں ہو رہیاں نیں۔

%%K t07-c2-04-2 p
نیویں کندھ اُتے اک سوہݨا \VUgras{مور} \textsuperscript{(18)} اپݨے شاندار پر وکھاؤݨ
نوں پُوچھ کھِلاردا اے۔

%%K t07-tit-7 sub
\VUpk{10.2pt}{40}{پیلی دا وِہڑا۔}
\VUsaut{3.0mm}

%%K t07-c2-05-1 p
5. --- \VUgras{طویلے} \textsuperscript{(19)} دے دروازے کھلے نیں، تے کھُرلی تے
\VUgras{وِچھاؤݨ دی پرالی} \textsuperscript{(23)} وکھائی دیندی اے اوس \VUgras{گھوڑی}
\textsuperscript{(21)} دی، جیہنوں \VUgras{سائیس} \textsuperscript{(20)} نے ہُݨے کھولیا
اے۔ \VUgras{وچھیرا} \textsuperscript{(22)}، اپݨیاں لمیاں لتاں اُتے اِنّا مزے دار،
چھالاں مار دا اپݨی ماں دے پچھے ٹُردا اے۔

%%K t07-c2-06-1 p
6. --- \VUgras{کھوہ} \textsuperscript{(29)} کول \VUgras{گاواں} \textsuperscript{(25)}
اوس لکڑ دی \VUgras{نانْد} \textsuperscript{(26)} توں پاݨی پیݨ جاندیاں نیں جیہڑی
اوہناں دا گھاٹ اے۔ \VUgras{وچھا} \textsuperscript{(24)} ایسے موقعے اُتے دُدھ پی
لیندا اے، تے \VUgras{بلد} \textsuperscript{(27)}، چارے دی چنگی وا لبھ کے، خوشی نال
اڑینگدا اے تے اپݨے تکڑے \VUgras{سِنگاں} \textsuperscript{(28)} والا سر مان نال
چُکدا اے۔

%%K t07-c2-07-1 p
7. --- سارے کم وچ لگے نیں۔ \VUgras{نوکراݨی} \textsuperscript{(32)} کھوہ دا
\VUgras{رسا} \textsuperscript{(31)} پھڑی ہوئی اے۔ اوہ ڈنگراں نوں پیاؤݨ لئی اک
\VUgras{بالٹی} \textsuperscript{(30)} وچ پاݨی کھِچ رہی اے۔ \VUgras{کوٹھے}
\textsuperscript{(33)} کول اک بندہ کھِلری پرالی کٹھی کر رہیا اے، تے
\VUgras{پیلی دے گھر} \textsuperscript{(34)} دے دروازے دے سامݨے اک بُڈھی
\VUgras{کتّݨ والی} \textsuperscript{(35)}، اپݨے چرخے تے اپݨی \VUgras{تکلے} نال،
پرانے دناں وانگوں \VUgras{سݨ} یا \VUgras{پٹسݨ} کت رہی اے۔

%%K t07-tit-8 sub
\VUsaut{3.0mm}
\VUpk{10.2pt}{40}{پیلی دا مالک۔}
\VUsaut{3.0mm}

%%K t07-c2-08-1 p
8. --- \VUgras{پیلی دا مالک} \textsuperscript{(36)} ایہہ سارا کم چلاندا اے تے اپݨے
بندیاں نوں حکم دیندا اے۔ اوہ اوہناں نوں آکھدا اے پئی \VUgras{سُہاگا}
\textsuperscript{(40)} تے \VUgras{کہی} \textsuperscript{(39)} ڈھک کے رکھ دیݨ، جیہڑے
اوس \VUgras{گھر} \textsuperscript{(38)} کول چھڈ دِتے گئے نیں جیہڑا \VUgras{کُتے}
\textsuperscript{(37)} دا اے۔

%%K t07-c2-09-1 p
9. --- اوہدی ہاک نال اک \VUgras{کبوتر} \textsuperscript{(41)} چونک اُٹھدا اے، جیہڑا
پوری چال نال \VUgras{کبوتر خانے} \textsuperscript{(42)} ول اُڈ جاندا اے، تے
\VUgras{کُکڑیاں} \textsuperscript{(43)} وی، جیہڑیاں چھیتی چھیتی \VUgras{ڈربے}
\textsuperscript{(44)} وچ مُڑ جاندیاں نیں۔

%%K t07-c2-10-1 p
10. --- \VUgras{بھیڈاں دے واڑے} \textsuperscript{(48)} دے سامݨے، ایہہ رہیا بُڈھا
\VUgras{آجڑی} \textsuperscript{(45)} اپݨا \VUgras{اِجّڑ} \textsuperscript{(49)}
موڑدا ہویا۔ اپݨی \VUgras{لاٹھی} \textsuperscript{(46)} پھڑی، جیہڑی اوہدا ساتھ اوہنا
ای نئیں چھڈدی جِنّا اوہدا فِکا \VUgras{کُڑتا} \textsuperscript{(47)}، اوہ واڑے ول
\VUgras{چھترے} \textsuperscript{(50)}، \VUgras{بھیڈاں} \textsuperscript{(53)}،
\VUgras{لیلے} \textsuperscript{(54)}، \VUgras{بکریاں} \textsuperscript{(51)} تے
\VUgras{بکری دے بچے} \textsuperscript{(52)} ہکدا اے۔ اوہدا وفادار \VUgras{کُتا}
\textsuperscript{(55)} آرگس سب توں پچھے آندا اے تے بھونک کے پچھڑیاں نوں اگے کردا اے۔

%%K t07-c2-11-1 p
11. --- ایہہ سارا رولا رپا اوس بھلی \VUgras{دُدھ والی گاں} \textsuperscript{(56)} دی
شانتی نئیں وگاڑدا، جیہڑی اپݨی \VUgras{ٹلی} \textsuperscript{(57)} وجاندی رہندی اے،
جدوں اوہنوں \VUgras{گؤ شالے} \textsuperscript{(58)} دے دروازے دے سامݨے چویا جاندا اے۔

%%K t07-c2-12-1 p
12. --- اوہ جِویں سمجھدی اے پئی اوہنے اپݨا دُدھ دیݨا ای اے، تاں جو
\VUgras{پیلی دی مالکݨ} \textsuperscript{(60)} \VUgras{مدھاݨی} \textsuperscript{(61)}
وچ \VUgras{مکھݨ} بݨا سکے، یا اوہ ودھیا \VUgras{پنیر} \textsuperscript{(64)} جیہڑے
\VUgras{ٹب} \textsuperscript{(63)} اُتے رکھے پٹرے اُتے نتھر رہے نیں۔

%%K t07-c2-13-1 p
13. --- سارا \VUgras{دُدھ گھر} \textsuperscript{(59)} \VUgras{پیلی دے مزدور}
\textsuperscript{(65)} دی بدولت بہت صاف رہندا اے، جیہنے ہُݨے اوس وڈی بالٹی وچ
\VUgras{دُدھ دے ڈبے} \textsuperscript{(62)} دھوتے نیں جیہڑی اوہ چُکی ہویا اے۔

%%K t07-tit-9 sub
\VUsaut{3.0mm}
\VUpk{10.2pt}{40}{کُکڑ خانہ۔}
\VUsaut{3.0mm}

%%K t07-c2-14-1 p
14. --- اپݨیاں بھارِیاں کھڑاواں وچ اوہ کجھ بھدے ڈھنگ نال ٹُردا اے، تے ایس نال
\VUgras{ٹرکی} \textsuperscript{(68)}، \VUgras{بطخاں} \textsuperscript{(66)}،
\VUgras{بطخاں دے نر} \textsuperscript{(67)} تے \VUgras{ہنس} \textsuperscript{(69)}
اُڈ جاندے نیں، جیہڑے اپݨیاں \VUgras{چُنجاں} \textsuperscript{(70)} چوڑیاں کھول کے بے
چین رولا پاندے نیں۔

%%K t07-c2-14-2 p
\VUgras{کُکڑ} \textsuperscript{(71)}، جیہڑا ودھ لڑاکا اے، اپݨی لال \VUgras{کلغی}
\textsuperscript{(72)} چُکدا اے، اپݨی \VUgras{پُوچھ} \textsuperscript{(73)} دے پر جھاڑدا
اے تے اک گونجدی « بانگ » چھڈدا اے۔

%%K t07-c2-15-1 p
15. --- اک \VUgras{کُکڑی ماں} \textsuperscript{(74)} \VUgras{گوہے دے ڈھیر}
\textsuperscript{(81)} اُتے چُگ رہی اے، اپݨے \VUgras{چوچیاں} \textsuperscript{(75)}
نال۔ \VUgras{سُور دا بچہ} \textsuperscript{(78)} اپݨی ماں ول دوڑ جاندا اے، اوس وڈی
\VUgras{سُوری} \textsuperscript{(77)} ول جیہڑی سانوں واڑے کول وکھائی دیندی اے۔

%%K t07-c2-16-1 p
16. --- اپݨے ڈربیاں وچ \VUgras{سہے} \textsuperscript{(79)} اوہ کُونا \VUgras{گھاہ}
\textsuperscript{(80)} چاء نال کھا رہے نیں جیہڑا مالک دی دھی اوہناں لئی لیاندی اے۔
جینی نوں اپݨے سہے بہت پیارے نیں، تے ہر دن، ڈربے توں تازے \VUgras{آنڈے}
\textsuperscript{(76)} لیاؤݨ مگروں، اوہ اپݨے نِکے یاراں نوں ملݨ آندی اے۔
\VUsaut{6.0mm}
\VUfilet{22mm}
